\bcpchaptersub{Solemnization of Matrimony}{The Form of Solemnization of Matrimony}

\begin{rubric}
First the banns of all that are to be married together must be published in the Church three several Sundays, during the time of Morning Service, or of Evening Service, (if there be no Morning Service,) immediately after the second Lesson; the Curate saying after the accustomed manner,
\end{rubric}

\lett{I}{ publish} the Banns of Marriage between \changeable{N.} of \changeable{---} and \changeable{N.} of \changeable{---}. If any of you know cause or just impediment, why these two persons should not be joined together in Holy Matrimony, ye are to declare it.
This is the first [second or third] time of asking.

\begin{rubric}
And if the persons that are to be married dwell in divers Parishes, the Banns must be asked in both Parishes; and the Curate of the one Parish shall not solemnize Metrimony betwixt them, without a Certificate of the Banns being thrice asked, from the Curate of the other Parish.
\end{rubric}

\begin{leftbar}
\begin{rubric}
Nor shall any be married without banns first asked, save in pursuance of a licence granted by ecclesiatical authority, or a certificate of due notice of the marriage, issued by a Superintendent Registrar; which licence or certificate, as the case may be, shall be produced to the Curate before the solemnization of the marriage.
\end{rubric}
\end{leftbar}

\begin{rubric}
It is convenient that the new-married persons should receive the holy Communion at the time of their Marriage, or at the first opportunity after their Marriage.
\end{rubric}

Dearly beloved, we are gathered together here in the sight of God, his angels and all the saints,
and in the face of the Church, to join together two bodies, to wit, those of this man and this
woman in holy Matrimony; which is an honourable estate, instituted of God in the time of man’s
innocency, signifying unto us the mystical union that is betwixt Christ and his Church; which holy
estate Christ adorned and beautified with his presence, and first miracle that he wrought, in Cana
of Galilee; and is commended of Saint Paul to be honourable among all men: and therefore is not
by any to be enterprised, nor taken in hand, unadvisedly, lightly, or wantonly, to satisfy men’s
carnal lusts and appetites, like brute beasts that have no understanding; but reverently, discreetly,
advisedly, soberly, and in the fear of God; duly considering the causes for which Matrimony was
ordained.\\
First, It was ordained for the procreation of children, to be brought up in the fear and nurture of
the Lord, and to the praise of his holy Name.\\
Secondly, It was ordained for a remedy against sin, and to avoid fornication; that such persons as
have not the gift of continency might marry, and keep themselves undefiled members of Christ’s
body.\\
Thirdly, It was ordained for the mutual society, help, and comfort, that the one ought to have of
the other, both in prosperity and adversity. Into which holy estate these two persons present
come now to be joined. Therefore if any man can shew any just cause, why they may not lawfully
be joined together, let him now speak, or else hereafter for ever hold his peace.

\begin{rubric}
And also, speaking unto the persons that shall be married, he shall say,
\end{rubric}

I require and charge you both, as ye will answer at the dreadful day of judgement, when the secrets of all hearts shall be disclosed, that if either of you know any impediment, why ye may not be lawfully joined together in Matrimony, ye do now confess it.
For be ye well assured, that so many as are coupled together otherwise than God's Word doth allow are not joined together by God; neither is their Matrimony lawful.

\begin{rubric}
At which day of Marriage, if any man do allege and declare any impediment, why they may not be coupled together in Matrimony, by God's law, or the laws of this Realm; and will be bound, and sufficient sureties with him, to the parties; or else put in a caution (to the full value of such charges as the persons to be married do thereby sustain) to prove his allegation: then the solemnization must be deferred, until such time as the truth be tried.
\end{rubric}

\begin{rubric}
If no impediment be alleged, then shall the Priest say unto the Man,
\end{rubric}

N. \lett{W}{ilt} thou have this woman to thy wedded wife, to live together after God's ordinance in the holy estate of Matrimony?
Wilt thou love her, comfort her, honour, and keep her in sickness and in health; and, forsaking all other, keep thee only unto her, so long as ye both shall live?

